\chapter{问题背景与研究目标}

\section{问题来源}
地下停车场发生火灾、烟雾扩散或结构异常时，救援队员和被困人员往往位于 GNSS 无法覆盖的区域。传统依赖卫星定位的方法在该场景失效，而仅依赖人工语音定位又存在延迟与误判。近年来，基于蓝牙信标（BLE Beacon）的室内定位方案部署成本较低，适合停车场这类规则空间，因此具有较高工程价值。

\section{研究难点}
在真实停车场环境中，定位算法面临以下挑战：
\begin{enumerate}
  \item 金属车体和立柱导致多路径反射，RSSI 波动大；
  \item 车辆动态进出造成遮挡变化，观测噪声呈非平稳特征；
  \item 单时刻点位估计抖动明显，难以形成可用轨迹；
  \item 应急场景要求算法实时、稳定、可解释。
\end{enumerate}

\section{问题定义}
设人员在二维平面上的位置为
\[
\bm p_t=[x_t,y_t]^{\T},\quad t=1,2,\dots,T.
\]
已知第 $i$ 个蓝牙信标坐标 $\bm b_i=[u_i,v_i]^{\T}$，在时刻 $t$ 观测到 RSSI 值 $z_{i,t}$。同时可通过步行模型给出短时运动先验（步长与方向变化约束）。目标是在噪声和缺失观测下估计整段轨迹 $\{\bm p_t\}_{t=1}^T$。

\section{本文目标}
本文目标包括：
\begin{itemize}
  \item 建立 RSSI 与位置的观测模型，并完成参数标定；
  \item 构建“观测项 + 运动约束项”的加权优化目标；
  \item 设计鲁棒求解算法，实现实时轨迹估计；
  \item 在不同噪声与遮挡条件下验证模型有效性。
\end{itemize}

\section{技术路线}
整体路线为“场景抽象 -- 数学建模 -- 鲁棒优化 -- 实验评估”。具体包括：
\begin{enumerate}
  \item 基于对数路径损耗模型把 RSSI 转为距离观测；
  \item 采用多信标加权最小二乘进行时刻级定位；
  \item 引入步长和转向约束构造滑动窗口平滑；
  \item 通过迭代重加权抑制粗差并输出稳定轨迹。
\end{enumerate}

\section{章节安排}
第 2 章给出观测模型与约束建模，第 3 章推导目标函数与优化量，第 4 章介绍初始化策略，第 5 章说明求解流程，第 6 至 8 章分别给出实验结果、不确定性分析和扩展讨论，第 9 章总结全文。
